\section{Weierstrass 第二逼近定理与 Vallee-Poussin 算子}

前面我们用 Bernstein 多项式证明了 Weierstrass 第一逼近定理，本节引入 Vallee-Poussin 算子，给出 Weierstrass 第二逼近定理的构造性证明。

\subsection{Vallee-Poussin 算子的定义与三角多项式表示}

定义 Vallee-Poussin 算子（简称 V-P 算子）
\[
V_n(f;x) = \frac{(2n)!!}{2\pi\cdot(2n-1)!!}\int_{-\pi}^{\pi} f(t)\cos^{2n}\frac{t-x}{2}\,dt,
\]
其中 $f$ 是以 $2\pi$ 为周期的连续函数。首先说明 $V_n(f;x)$ 是一个三角多项式。

利用 Euler 公式
\[
\cos\frac{t-x}{2} = \frac{e^{i\frac{t-x}{2}} + e^{-i\frac{t-x}{2}}}{2},
\]
由二项式定理，
\[
\cos^{2n}\frac{t-x}{2} = \frac{1}{2^{2n}} \sum_{k=0}^{2n} \binom{2n}{k} \left(e^{i\frac{t-x}{2}}\right)^{2n-k} \left(e^{-i\frac{t-x}{2}}\right)^k.
\]
代入积分表达式，记 $C_n = \frac{(2n)!!}{2\pi\cdot(2n-1)!!}$，得
\[
V_n(f;x) = \frac{C_n}{2^{2n}} \sum_{k=0}^{2n} \binom{2n}{k} \int_{-\pi}^{\pi} f(t)\, e^{i(n-k)(t-x)}\,dt.
\]
由于是有限项求和，积分与求和可交换次序，故
\[
V_n(f;x) = \sum_{k=0}^{2n} \left[ \frac{C_n \binom{2n}{k}}{2^{2n}} \int_{-\pi}^{\pi} f(t)\, e^{i(n-k)t}\,dt \right] e^{-i(n-k)x}.
\]
令 $j = n-k$，则求和范围变为 $j = -n, \dots, n$，即
\[
V_n(f;x) = \sum_{j=-n}^{n} \left[ \frac{C_n \binom{2n}{n-j}}{2^{2n}} \int_{-\pi}^{\pi} f(t)\, e^{ijt}\,dt \right] e^{-ijx}.
\]
括号内的积分与 $x$ 无关，记
\[
c_j = \frac{C_n \binom{2n}{n-j}}{2^{2n}} \int_{-\pi}^{\pi} f(t)\, e^{ijt}\,dt,
\]
则 $V_n(f;x) = \sum_{j=-n}^{n} c_j e^{-ijx}$。

注意到
\[
c_{-j} = \frac{C_n \binom{2n}{n+j}}{2^{2n}} \int_{-\pi}^{\pi} f(t)\, e^{-ijt}\,dt = \overline{c_j},
\]
因此系数满足共轭对称性。从而
\[
c_j e^{-ijx} + c_{-j} e^{ijx} = c_j e^{-ijx} + \overline{c_j} e^{ijx} = 2\,\text{Re}(c_j e^{-ijx}),
\]
这表明 $V_n(f;x)$ 可以完全用实的三角多项式表示。因此 V-P 算子确实构造出了一列三角多项式。

\subsection{一致收敛性的证明}

由 Wallis 公式，
\begin{align*}
\int_{-\pi}^{\pi} \cos^{2n}\frac{t-x}{2}\,dt
&= \int_{-\pi}^{\pi} \cos^{2n}\frac{t}{2}\,dt
= 2\int_{0}^{\pi} \cos^{2n}\frac{t}{2}\,dt \\
&= 4\int_{0}^{\frac{\pi}{2}} \cos^{2n} u\,du
= \frac{2\pi\cdot(2n-1)!!}{(2n)!!},
\end{align*}
恰好是 $V_n(f;x)$ 前面系数的倒数。由此可将 $f(x)$ 表示为
\[
f(x) = \frac{(2n)!!}{2\pi\cdot(2n-1)!!}\int_{-\pi}^{\pi} f(x)\cos^{2n}\frac{t-x}{2}\,dt.
\]

作差并取绝对值，
\begin{align*}
|V_n(f;x) - f(x)|
&= \left| \frac{(2n)!!}{2\pi\cdot(2n-1)!!}\int_{-\pi}^{\pi} (f(t)-f(x))\cos^{2n}\frac{t-x}{2}\,dt \right| \\
&\leq \frac{(2n)!!}{2\pi\cdot(2n-1)!!}\int_{-\pi}^{\pi} |f(t)-f(x)|\cos^{2n}\frac{t-x}{2}\,dt.
\end{align*}

由于 $f$ 是 $\mathbb{R}$ 上以 $2\pi$ 为周期的连续函数，必一致连续：$\forall\epsilon>0,\exists\delta>0$，当 $|x-y|<\delta$ 时，$|f(x)-f(y)|<\epsilon$。

将积分区域划分为两部分。令
\[
I_1 = [-\pi,\pi] \cap [x-\delta, x+\delta],\quad
I_2 = [-\pi,\pi] \setminus I_1,
\]
并记 $I = I_1 \cup I_2 = [-\pi,\pi]$。

\textbf{对 $I_1$ 的估计}：在 $I_1$ 上 $|t-x|<\delta$，故 $|f(t)-f(x)|<\epsilon$，于是
\begin{align*}
\frac{(2n)!!}{2\pi\cdot(2n-1)!!}\int_{I_1} |f(t)-f(x)|\cos^{2n}\frac{t-x}{2}\,dt
&\leq \frac{(2n)!!}{2\pi\cdot(2n-1)!!}\int_{I_1} \epsilon\cos^{2n}\frac{t-x}{2}\,dt \\
&\leq \epsilon \cdot \frac{(2n)!!}{2\pi\cdot(2n-1)!!}\int_{I}\cos^{2n}\frac{t-x}{2}\,dt = \epsilon.
\end{align*}

\textbf{对 $I_2$ 的估计}：在 $I_2$ 上 $|t-x|\geq\delta$，因此 $\left|\cos\frac{t-x}{2}\right| \leq \left|\cos\frac{\delta}{2}\right|$。令 $M = \sup_{x\in\mathbb{R}}|f(x)|$，则 $|f(t)-f(x)|\leq 2M$。于是
\begin{align*}
&\frac{(2n)!!}{2\pi\cdot(2n-1)!!}\int_{I_2} |f(t)-f(x)|\cos^{2n}\frac{t-x}{2}\,dt \\
&\leq \frac{(2n)!!}{2\pi\cdot(2n-1)!!}\int_{I} 2M\cos^{2n}\frac{\delta}{2}\,dt \\
&= \frac{2M\cdot(2n)!!\cdot\cos^{2n}\frac{\delta}{2}\cdot 2\pi}{2\pi\cdot(2n-1)!!}
= \frac{2M\cdot(2n)!!\cdot\cos^{2n}\frac{\delta}{2}}{(2n-1)!!}.
\end{align*}

利用归纳法可证 $\frac{(2n)!!}{(2n-1)!!} \leq 2n$，因此上述积分不超过 $4Mn\cos^{2n}\frac{\delta}{2}$。

由于 $\cos^{2n}\frac{\delta}{2} < 1$，当 $n\to\infty$ 时该项趋于 $0$。故存在 $N$，使得当 $n>N$ 时，该部分小于 $\epsilon$。

综合两部分，当 $n>N$ 时，
\[
|V_n(f;x) - f(x)| < 2\epsilon
\]
对任意 $x\in\mathbb{R}$ 成立。这就证明了 $V_n(f;x)$ 一致收敛于 $f(x)$，从而完成了 Weierstrass 第二逼近定理的构造性证明。

\begin{remark}
还可以用 Weierstrass 第一逼近定理来推导第二逼近定理，方法不止一种。
\end{remark}